%% kapa: exam=1 round=19 year=2008 period=1 form=A subject=civil_law subject_ko=민법 seq=27 number=27
%% source: Q-Net 감정평가사 기출문제 (cst00309) artlSeq=1001962 file=extracted/1008/1001962/2008년 제19회 감정평가사 1차시험 기출문제/1교시.hwp
%% conversion: kordoc 4.24.51 -> tools/gen1.py (auto); review=proofread-sonnet; flags=none
\begin{question}{27}
물권의 소멸에 관한 설명으로 틀린 것은? (다툼이 있는 경우에는 판례에 의함)
\begin{choices}
  \item 유치물에 관하여 이해관계를 가지고 있는 자인 채무자나 유치물의 소유자가 상당한 담보를 제공하면 유치권 소멸청구의 의사표시를 할 수 있다.
  \item 토지가 포락되어 그 효용을 상실한 경우에는 그 후 포락된 토지가 성토되더라도 종전의 소유자가 토지소유권을 다시 취득할 수 없다.
  \item 승역지 소유자가 그 소유권을 지역권자에게 이전시키는 의사표시를 하고 등기를 경료한 경우에 지역권은 소멸된다.
  \item 근저당권자가 소유권을 취득하면 근저당권은 혼동에 의하여 절대적으로 소멸하므로, 그 소유권취득이 무효이더라도 소멸하였던 근저당권은 부활하지 않는다.
  \item 대항력을 갖춘 임대차의 목적물에 저당권이 설정된 후, 임차인이 임차목적물의 소유권을 취득하더라도 임차권은 혼동으로 소멸되지 않는다.
\end{choices}
\end{question}
\answer{4}
